% ==================== RÉSUMÉ ARABE ====================
\chapter*{Résumé (Arabe)}
\begin{center}
\LARGE \textarabic{ملخص}
\end{center}
\addcontentsline{toc}{chapter}{Résumé (Arabe)}
\thispagestyle{plain}

\begin{center}
\begin{minipage}{0.92\textwidth}
\begin{Arabic}
تُقدّم هذه الأطروحة تصميم وتطوير منصة "روح وين بجيت"، وهي منصة ذكية لمشاركة السيارات عبر الهاتف المحمول، مصممة لمعالجة تحديات النقل بين الولايات في الجزائر. يستجيب هذا العمل لحاجة مجتمعية ملحة ناتجة عن غياب حل رقمي آمن وموثوق به للسفر لمسافات طويلة في جميع أنحاء البلاد. غالبًا ما تتسم خيارات النقل الحالية بارتفاع التكاليف، ومحدودية التوافر، وعدم كفاية آليات بناء الثقة بين الركاب والسائقين. تهدف المنصة المقترحة إلى تحسين إمكانية الوصول إلى خدمات التنقل المشترك، وسلامتها، وكفاءتها، من خلال دمج ميزات تُسهّل تخطيط الرحلات، وحجز الرحلات، والتواصل الفوري، ومراقبة الرحلات. يُولى اهتمام خاص لثقة المستخدم وأمانه من خلال تطبيق آليات التحقق من الهوية، وتدابير منع الاحتيال، وإدارة المعاملات بشكل آمن. كما تتضمن المنصة وظائف ذكية لدعم اتخاذ القرارات وتحسين تجربة المستخدم بشكل عام. تم تصميم وتنفيذ بنية نظام شاملة لضمان قابلية التوسع، والموثوقية، والأداء. يدعم الحل المُطوّر التفاعلات الفورية بين المستخدمين، والخدمات القائمة على الموقع، وإدارة الرحلات الديناميكية، وخيارات دفع متعددة مُكيّفة مع السياق الجزائري. تم تقييم فعالية المنصة من خلال اختبارات شاملة، تضمنت التحقق الآلي من الوظائف الأساسية، واختبار التكامل، والتحقق من واجهة المستخدم. تُظهر النتائج المُتحصل عليها فعالية الحل المقترح وقدرته على تلبية المتطلبات الوظيفية والأمنية لخدمة مشاركة السيارات الحديثة. إجمالاً، يُسهم هذا العمل في التحول الرقمي لخدمات النقل في الجزائر من خلال اقتراح حل عملي ومبتكر للتنقل يُعزز السفر بين الولايات بشكل أكثر أمانًا وسهولة واستدامة.

\vspace{0.5cm}
\textbf{الكلمات المفتاحية~:} مشاركة السيارات، التنقل الذكي، تطبيقات الهاتف المحمول، أنظمة النقل، ثقة المستخدم، التحقق من الهوية، الذكاء الاصطناعي، المدفوعات الرقمية، الجزائر، النقل بين الولايات.
\end{Arabic}
\end{minipage}
\end{center}

% ==================== RÉSUMÉ FRANÇAIS ====================
\chapter*{Résumé}
\addcontentsline{toc}{chapter}{Résumé}
\thispagestyle{plain}

Cette thèse présente la conception et le développement de RohWinBghit, une plateforme mobile intelligente de covoiturage conçue pour répondre aux enjeux du transport inter-wilayas en Algérie. Ce travail répond à un besoin sociétal majeur né de l'absence d'une solution numérique sécurisée, fiable et digne de confiance pour les longs trajets à travers le pays. Les options de transport existantes sont souvent caractérisées par des coûts élevés, une disponibilité limitée et des mécanismes insuffisants pour instaurer la confiance entre passagers et conducteurs. La plateforme proposée vise à améliorer l'accessibilité, la sécurité et l'efficacité des services de mobilité partagée en intégrant des fonctionnalités facilitant la planification des trajets, la réservation de courses, la communication en temps réel et le suivi des déplacements. Une attention particulière est portée à la confiance et à la sécurité des utilisateurs grâce à la mise en œuvre de mécanismes de vérification d'identité, de mesures de prévention de la fraude et d'une gestion sécurisée des transactions. La plateforme intègre également des fonctionnalités intelligentes pour faciliter la prise de décision et améliorer l'expérience utilisateur globale. Une architecture système complète a été conçue et mise en œuvre pour garantir l'évolutivité, la fiabilité et la performance. La solution développée prend en charge les interactions en temps réel entre les utilisateurs, les services de géolocalisation, la gestion dynamique des trajets et de multiples options de paiement adaptées au contexte algérien. L'efficacité de la plateforme a été évaluée par des tests approfondis, incluant la validation automatisée des fonctionnalités essentielles, des tests d'intégration et la vérification de l'interface utilisateur. Les résultats obtenus démontrent la robustesse de la solution proposée et sa capacité à répondre aux exigences fonctionnelles et de sécurité d'un service de covoiturage moderne. Ce travail contribue à la transformation numérique des services de transport en Algérie en proposant une solution de mobilité pratique et innovante qui favorise des déplacements inter-wilayas plus sûrs, plus accessibles et plus durables.

\vspace{0.3cm}
\textbf{Mots-clés~:} Covoiturage, Mobilité intelligente, Application mobile, Systèmes de transport, Confiance des utilisateurs, Vérification d'identité, Intelligence artificielle, Paiements numériques, Algérie, Transport inter-wilayas.